\documentclass[letterpaper, 10pt]{article}
\usepackage{graphicx} % Required for inserting images
\usepackage[hidelinks]{hyperref} % Required for hyperlinks
\usepackage[top=0.5in, bottom=0.5in, left=0.7in, right=0.7in]{geometry} % Page margins
\pagestyle{empty} % No page numbers
\begin{document}
%------------HEADING------------
\begin{center}
\textbf{\LARGE Chung MinKi} \\ \vspace{2pt}
\href{mailto:spacemoon@skku.edu}{\underline{spacemoon@skku.edu}} $|$
\href{https://github.com/justinkays}{\underline{github.com}}
\end{center}
%-----------TECHNICAL SKILLS-----------
\section*{Technical Skills}
\hrule
\vspace{4pt}
\begin{itemize}
% \item[] : same as \item but with no bullet point ([] replaces the bullet symbol with nothing)
\item[] \textbf{Programming}: C++, Python% List your programming languages
\item[] \textbf{Frameworks \& Tools}: PyTorch, OpenCV, Arduino, Git, Docker% List your frameworks and tools
\item[] \textbf{Languages}: Korean, English% List your spoken languages
\end{itemize}
%------------EXPERIENCE------------
\section*{Experience}
\hrule
\vspace{4pt}
\begin{itemize}
\item[]
% \hfill : fills the remaining horizontal space, pushing text to the right
\textbf{Undergraduate Researcher (Industry-Academic Cooperation Project)} \hfill 2026 \\
\textit{Sungkyunkwan University (SKKU)} \hfill \textit{Suwon, South Korea}
\begin{itemize}
\item Researched frequency-spectrum approaches for detecting AI-generated visual content using Python, PyTorch, and OpenCV.
\item Published and presented ``An Empirical Study of Frequency Spectrum for Detection of AI Generated Visual Content'' at the CISC conference.
\end{itemize}
\end{itemize}
%-----------SELECTED PROJECTS-----------
% Include: project name, brief description, and your role
\section*{Selected Projects}
\hrule
\vspace{4pt}
\begin{itemize}
\item[]
\textbf{Frequency Spectrum Analysis for AI-Generated Visual Content Detection} \hfill October 2026 \\
\textit{Role: Researcher \& Presenter} \hfill \textit{Yeosu, South Korea}
\begin{itemize}
\item Built an empirical evaluation workflow to analyze frequency-spectrum characteristics of AI-generated visual content.
\item Authored and presented the resulting study at the Conference on Information Security and Cryptology (CISC).
\end{itemize}
\item[]
\textbf{Autonomous Maze-Escape Robot} \hfill March 2022 -- June 2022 \\
\textit{Role: Team Developer (Creative Engineering Design Competition)} \hfill \textit{OpenCV, Arduino}
\begin{itemize}
\item Built an autonomous maze-escaping robot combining OpenCV-based vision processing with Arduino hardware control.
\item Won 2nd Place in the Creative Engineering Design Competition.
\end{itemize}
\end{itemize}
%-----------HONORS & AWARDS-----------
% Include: awards, certifications, or honors you have
\section*{Honors \& Awards}
\hrule
\vspace{4pt}
\begin{itemize}
\item \textbf{2nd Place, Creative Engineering Design Competition} -- Autonomous Maze-Escape Robot (OpenCV, Arduino) \hfill 2022
\end{itemize}
%------------EDUCATION------------
% Include: university name, degree, major, relevant courses, academic accomplishments, etc
\section*{Education}
\hrule
\vspace{4pt}
\begin{itemize}
\item[]
\textbf{Sungkyunkwan University (SKKU)} \hfill February 2026 -- Present \\
\textit{B.S. in Computer Science (Undergraduate)} \hfill \textit{Suwon, South Korea}
\end{itemize}
%-----------------------------------
\end{document}
